\section{Give each index one job}
The transpose chapter moved values from one address to another. Multiplication
combines values. The row index \(i\) chooses a row of \(A\), the column index
\(j\) chooses a column of \(B\), and the reduction index \(k\) selects the pair
being multiplied. The reduction consumes \(K\) products to form one output value.
It does not advance to another output until that sum is complete in the baseline.

For an \(M\)-by-\(K\) left input, a complete row contains \(K\) words. For a
\(K\)-by-\(N\) right input, a complete row contains \(N\) words. These widths
need not match. A square example can conceal a mistaken stride, because substituting
one dimension for another leaves the address unchanged when all three are equal.
Use a rectangular example before timing any version.

\begin{MARSWorkedExample}{Two rows, three products, two output columns}
\[
 A=\begin{pmatrix}1&2&3\\4&5&6\end{pmatrix},\qquad
 B=\begin{pmatrix}7&8\\9&10\\11&12\end{pmatrix}.
\]
For output row zero, column one, use the first row of \(A\) and the second
column of \(B\). Their products are \(1\cdot8\), \(2\cdot10\) and
\(3\cdot12\), so \(C_{0,1}=64\). The complete result is
\[
 C=\begin{pmatrix}58&64\\139&154\end{pmatrix}.
\]
Flattened row-major output is therefore \(58,64,139,154\), with eight bytes
between neighboring elements. The right input remains in its original row-major
layout; reading a column does not require first transposing the array.
\end{MARSWorkedExample}

\begin{MARSExercise}{Predict one cell before running}
Write out the three products that form \(C_{1,0}\). Record the input indices and
values separately. Then explain why multiplying \(A_{1,k}\) by \(B_{0,k}\)
would select a row of \(B\) and violate this product's definition.
\end{MARSExercise}
